    \begin{subsection}{Properties, Examples and Obstructions}
    Risk measures are now a standard tool for evaluation in actuarial science as well as in financial mathematics as a whole. 
    One of the most basic risk measures,
    \[ 
        \rho(X) = \mathbb{E}[X] 
    \] 
    dates back as early as 1709 to Nicolaus I Bernoulli’s PhD. thesis \cite{bernoulli1709dissertatio}, in which he proposed it as the fair price for insurance contracts.
    % https://pmc.ncbi.nlm.nih.gov/articles/PMC11928225/
    % Roman Law on the Just Price in Nicolaus Bernoulli’s Mathematics
    Remarkably, only a few years later, Bernoulli also suggested an adjustment based on individual preferences—a concept we would now recognize as a precursor to utility theory.
    It took several centuries, however, until Markowitz’s portfolio theory \cite{markowitzportfolio}, which introduced variance as a measure of risk, reignited the theory and 
    led to the early development of downside risk measures such as lower partial moments. Specifically, in 1959, Markowitz (Chapter 9 \cite{markowitzsemivariance}) proposed
    \[
        \mathbb{E}[(|\mathbb{E}[X] - X|\mathbf{1}_{ \{\mathbb{E}[X] - X > 0 \} })^k] 
    \] 
    Over the following decades, various risk measures flourished, mostly under the term "premium principle", 
    % A General Class of Three-Parameter Risk Measures Bernell K. Stone
    of which special mentioning deserves the value-at-risk 
    \[ 
        \rho(X) = \text{VaR}_{\lambda}(X) = \inf \{ m \in \mathbb{R} | \pp{X + m < 0} \leq \lambda \}
    \]
    and the expected shortfall, or average value at risk ,
    \[
        \rho(X) = \text{AVR}_{\lambda}(X) = \frac{1}{\lambda} \int_{0}^\lambda \text{VaR}_{\alpha}(X) d\alpha
    \]
    until in 2011 Föllmer and Schied \cite{foellmerschied} put the theory on rigorous axiomatic foundations.
    % Risk Measures and Theories of Choice Tsanakas
    While an algebraic approach to portfolio theory is already well developed—since it is inherently an optimization problem on an intersection of hyperplanes and a conic—
    the algebraic aspect of risk measures still offers ample room for development, to the best of the author’s knowledge. 
    \newline

    We first give some basic definitions for risk measures.
    Let \( \mathcal{X} = \mathcal{L}^p(\Omega, \mathcal{F}, \mathbb{P}) \) or subspace thereof, 
    for some suitable choice of \( p \).
    \begin{definition} \label{defriskmeasure}
    A mapping \( \rho : \mathcal{X} \to \mathbb{R} \) is called a monetary risk measure if it satisfies
    the following conditions for all \( X, Y \in \mathcal{X} \).
        \begin{enumerate}
            \item Monotonicity/Antitone: If \( X \leq Y\), then \(\rho(X) \geq \rho(Y) \).
            \item Cash invariance: If \( m \in \mathbb{R} \), then \( \rho(X + m) = \rho(X) - m \).
        \end{enumerate}
        It is called normalized if \( \rho(0) = 0 \), convex if \( \rho \) is a convex functional, and coherent if it is in addition positively homogenous,
        that is \( \rho(\lambda X) = \lambda \rho(X),\ \lambda \in \mathbb{R}_{\geq 0} \).
    \end{definition}
    \begin{remark}
    If we say risk measure, we mean monetary risk measure, if not otherwise stated. The concept of coherence could be referred to as a sublinear 
    risk measure in functional analytic terminology.
    \end{remark}
    \begin{definition}
        A risk measure \( \rho \) is
        \begin{enumerate}
            \item sensitive, if \( \rho(-X) > \rho(0) \) for \( X \in L^\infty_{+}(\Omega, \mathcal{F}, \mathbb{P}) \) and \( \pp{X > 0} > 0 \),
            \item law-invariant, if \( \rho(X) = \rho(Y) \) for \( X \sim Y\) under \( \mathbb{P} \),
            \item continuous from below, if \( X_n \nearrow X \implies \rho(X_n) \searrow \rho(X) \),
            \item continuous from above, if \( X_n \searrow X \implies \rho(X_n) \nearrow \rho(X) \).
        \end{enumerate}
    \end{definition}

    \begin{definition}[Linear cones]
    A set \( S \) in a real vector space \( \mathcal{X} \) is a cone if \( x \in S \implies \lambda x \in S \) for \( \lambda \in \mathbb{R}_{>0} \).
    \end{definition}
    \begin{remark}
       To differentiate cones as in the previous from the ones such as in Definition \ref{coneofrings}
       we call the first linear cones and the latter algebraic cones.
    \end{remark}
    \begin{definition}
        The acceptance set of a risk measure \( \rho \) is 
        \[ \mathcal{A}_\rho = \{ X \in \mathcal{X} | \rho(X) \leq 0 \}. \]
    \end{definition}
    \begin{proposition} \label{acceptancesetpropisriskprop}
        Assume \( \mathcal{A}_{\rho} \neq \emptyset \), it holds:
        \begin{enumerate}
            \item \( \rho(X) := \inf \{ m \in \mathbb{R} | m + X \in \mathcal{A}_{\rho} \}   \) 
            \item \( \rho \) is a convex risk measure if and only if \( \mathcal{A}_{\rho} \) is convex.
            \item \( \rho \) is positive homogenous if and only if \( \mathcal{A}_{\rho} \) is a cone.
            \item \( \rho \) is coherent if and only if \( \mathcal{A}_{\rho} \) is a convex cone 
        \end{enumerate}
    \end{proposition}
    \begin{proof}
        \cite{foellmerschied} Proposition 4.6
    \end{proof}
    \end{subsection}
